% Shared layout for marti_ferran_gonzales_cv.tex and marti_ferran_gonzales_cover_letter.tex.
% The CV defines \cvsetupfont before % Shared layout for marti_ferran_gonzales_cv.tex and marti_ferran_gonzales_cover_letter.tex.
% The CV defines \cvsetupfont before % Shared layout for marti_ferran_gonzales_cv.tex and marti_ferran_gonzales_cover_letter.tex.
% The CV defines \cvsetupfont before % Shared layout for marti_ferran_gonzales_cv.tex and marti_ferran_gonzales_cover_letter.tex.
% The CV defines \cvsetupfont before \input{preamble}.
% Author : Jake Gutierrez
% Based off of: https://github.com/sb2nov/resume
% License : MIT
% This file is not a document. Compile marti_ferran_gonzales_cv.tex or marti_ferran_gonzales_cover_letter.tex.

\ifdefined\cvsetupfont
  \usepackage{fontspec}
  \cvsetupfont
  \renewcommand{\familydefault}{\sfdefault}
\else
  \usepackage[T1]{fontenc}
  \usepackage[utf8]{inputenc}
\fi
\usepackage{latexsym}
\usepackage[empty]{fullpage}
\usepackage{titlesec}
\usepackage{marvosym}
\usepackage[usenames,dvipsnames]{color}
\usepackage{verbatim}
\usepackage{enumitem}
\usepackage[hidelinks]{hyperref}
\usepackage{fancyhdr}
\usepackage[english]{babel}
\usepackage{tabularx}
\ifdefined\cvsetupfont
\else
  \input{glyphtounicode}
\fi

\pagestyle{fancy}
\fancyhf{} % clear all header and footer fields
\fancyfoot{}
\renewcommand{\headrulewidth}{0pt}
\renewcommand{\footrulewidth}{0pt}

% Adjust margins
\addtolength{\oddsidemargin}{-0.5in}
\addtolength{\evensidemargin}{-0.5in}
\addtolength{\textwidth}{1in}
\addtolength{\topmargin}{-.5in}
\addtolength{\textheight}{1.0in}

\urlstyle{same}

\raggedbottom
\raggedright
\setlength{\tabcolsep}{0in}

% Sans faces here have no small caps, so the section title stays in the same font.
\ifdefined\cvsetupfont
  \titleformat{\section}{
    \vspace{-4pt}\raggedright\large
  }{}{0em}{\MakeUppercase}[\color{black}\titlerule \vspace{-5pt}]
\else
  \titleformat{\section}{
    \vspace{-4pt}\scshape\raggedright\large
  }{}{0em}{}[\color{black}\titlerule \vspace{-5pt}]
  % Ensure that the generated pdf is machine readable/ATS parsable
  \pdfgentounicode=1
\fi

%-------------------------
% Custom commands
\newcommand{\resumeItem}[1]{
  \item\small{
    {#1 \vspace{-2pt}}
  }
}

\newcommand{\resumeSubheading}[4]{
  \vspace{-2pt}\item
    \begin{tabular*}{0.97\textwidth}[t]{l@{\extracolsep{\fill}}r}
      \textbf{#1} & #2 \\
      \textit{\small#3} & \textit{\small #4} \\
    \end{tabular*}\vspace{-7pt}
}

\newcommand{\resumeSubSubheading}[2]{
    \item
    \begin{tabular*}{0.97\textwidth}{l@{\extracolsep{\fill}}r}
      \textit{\small#1} & \textit{\small #2} \\
    \end{tabular*}\vspace{-7pt}
}

\newcommand{\resumeProjectHeading}[2]{
    \item
    \begin{tabular*}{0.97\textwidth}{l@{\extracolsep{\fill}}r}
      \small#1 & #2 \\
    \end{tabular*}\vspace{-7pt}
}

\newcommand{\resumeSubItem}[1]{\resumeItem{#1}\vspace{-4pt}}

\renewcommand\labelitemii{$\vcenter{\hbox{\tiny$\bullet$}}$}

\newcommand{\resumeSubHeadingListStart}{\begin{itemize}[leftmargin=0.15in, label={}]}
\newcommand{\resumeSubHeadingListEnd}{\end{itemize}}
\newcommand{\resumeItemListStart}{\begin{itemize}}
\newcommand{\resumeItemListEnd}{\end{itemize}\vspace{-5pt}}

% Role context under a \resumeSubheading. The heading table is top-aligned and
% pulls the next line up, so the blurb must start a new paragraph below it.
\newcommand{\resumeIntro}[1]{\par\vspace{3pt}{\small #1\par}}
.
% Author : Jake Gutierrez
% Based off of: https://github.com/sb2nov/resume
% License : MIT
% This file is not a document. Compile marti_ferran_gonzales_cv.tex or marti_ferran_gonzales_cover_letter.tex.

\ifdefined\cvsetupfont
  \usepackage{fontspec}
  \cvsetupfont
  \renewcommand{\familydefault}{\sfdefault}
\else
  \usepackage[T1]{fontenc}
  \usepackage[utf8]{inputenc}
\fi
\usepackage{latexsym}
\usepackage[empty]{fullpage}
\usepackage{titlesec}
\usepackage{marvosym}
\usepackage[usenames,dvipsnames]{color}
\usepackage{verbatim}
\usepackage{enumitem}
\usepackage[hidelinks]{hyperref}
\usepackage{fancyhdr}
\usepackage[english]{babel}
\usepackage{tabularx}
\ifdefined\cvsetupfont
\else
  \input{glyphtounicode}
\fi

\pagestyle{fancy}
\fancyhf{} % clear all header and footer fields
\fancyfoot{}
\renewcommand{\headrulewidth}{0pt}
\renewcommand{\footrulewidth}{0pt}

% Adjust margins
\addtolength{\oddsidemargin}{-0.5in}
\addtolength{\evensidemargin}{-0.5in}
\addtolength{\textwidth}{1in}
\addtolength{\topmargin}{-.5in}
\addtolength{\textheight}{1.0in}

\urlstyle{same}

\raggedbottom
\raggedright
\setlength{\tabcolsep}{0in}

% Sans faces here have no small caps, so the section title stays in the same font.
\ifdefined\cvsetupfont
  \titleformat{\section}{
    \vspace{-4pt}\raggedright\large
  }{}{0em}{\MakeUppercase}[\color{black}\titlerule \vspace{-5pt}]
\else
  \titleformat{\section}{
    \vspace{-4pt}\scshape\raggedright\large
  }{}{0em}{}[\color{black}\titlerule \vspace{-5pt}]
  % Ensure that the generated pdf is machine readable/ATS parsable
  \pdfgentounicode=1
\fi

%-------------------------
% Custom commands
\newcommand{\resumeItem}[1]{
  \item\small{
    {#1 \vspace{-2pt}}
  }
}

\newcommand{\resumeSubheading}[4]{
  \vspace{-2pt}\item
    \begin{tabular*}{0.97\textwidth}[t]{l@{\extracolsep{\fill}}r}
      \textbf{#1} & #2 \\
      \textit{\small#3} & \textit{\small #4} \\
    \end{tabular*}\vspace{-7pt}
}

\newcommand{\resumeSubSubheading}[2]{
    \item
    \begin{tabular*}{0.97\textwidth}{l@{\extracolsep{\fill}}r}
      \textit{\small#1} & \textit{\small #2} \\
    \end{tabular*}\vspace{-7pt}
}

\newcommand{\resumeProjectHeading}[2]{
    \item
    \begin{tabular*}{0.97\textwidth}{l@{\extracolsep{\fill}}r}
      \small#1 & #2 \\
    \end{tabular*}\vspace{-7pt}
}

\newcommand{\resumeSubItem}[1]{\resumeItem{#1}\vspace{-4pt}}

\renewcommand\labelitemii{$\vcenter{\hbox{\tiny$\bullet$}}$}

\newcommand{\resumeSubHeadingListStart}{\begin{itemize}[leftmargin=0.15in, label={}]}
\newcommand{\resumeSubHeadingListEnd}{\end{itemize}}
\newcommand{\resumeItemListStart}{\begin{itemize}}
\newcommand{\resumeItemListEnd}{\end{itemize}\vspace{-5pt}}

% Role context under a \resumeSubheading. The heading table is top-aligned and
% pulls the next line up, so the blurb must start a new paragraph below it.
\newcommand{\resumeIntro}[1]{\par\vspace{3pt}{\small #1\par}}
.
% Author : Jake Gutierrez
% Based off of: https://github.com/sb2nov/resume
% License : MIT
% This file is not a document. Compile marti_ferran_gonzales_cv.tex or marti_ferran_gonzales_cover_letter.tex.

\ifdefined\cvsetupfont
  \usepackage{fontspec}
  \cvsetupfont
  \renewcommand{\familydefault}{\sfdefault}
\else
  \usepackage[T1]{fontenc}
  \usepackage[utf8]{inputenc}
\fi
\usepackage{latexsym}
\usepackage[empty]{fullpage}
\usepackage{titlesec}
\usepackage{marvosym}
\usepackage[usenames,dvipsnames]{color}
\usepackage{verbatim}
\usepackage{enumitem}
\usepackage[hidelinks]{hyperref}
\usepackage{fancyhdr}
\usepackage[english]{babel}
\usepackage{tabularx}
\ifdefined\cvsetupfont
\else
  \input{glyphtounicode}
\fi

\pagestyle{fancy}
\fancyhf{} % clear all header and footer fields
\fancyfoot{}
\renewcommand{\headrulewidth}{0pt}
\renewcommand{\footrulewidth}{0pt}

% Adjust margins
\addtolength{\oddsidemargin}{-0.5in}
\addtolength{\evensidemargin}{-0.5in}
\addtolength{\textwidth}{1in}
\addtolength{\topmargin}{-.5in}
\addtolength{\textheight}{1.0in}

\urlstyle{same}

\raggedbottom
\raggedright
\setlength{\tabcolsep}{0in}

% Sans faces here have no small caps, so the section title stays in the same font.
\ifdefined\cvsetupfont
  \titleformat{\section}{
    \vspace{-4pt}\raggedright\large
  }{}{0em}{\MakeUppercase}[\color{black}\titlerule \vspace{-5pt}]
\else
  \titleformat{\section}{
    \vspace{-4pt}\scshape\raggedright\large
  }{}{0em}{}[\color{black}\titlerule \vspace{-5pt}]
  % Ensure that the generated pdf is machine readable/ATS parsable
  \pdfgentounicode=1
\fi

%-------------------------
% Custom commands
\newcommand{\resumeItem}[1]{
  \item\small{
    {#1 \vspace{-2pt}}
  }
}

\newcommand{\resumeSubheading}[4]{
  \vspace{-2pt}\item
    \begin{tabular*}{0.97\textwidth}[t]{l@{\extracolsep{\fill}}r}
      \textbf{#1} & #2 \\
      \textit{\small#3} & \textit{\small #4} \\
    \end{tabular*}\vspace{-7pt}
}

\newcommand{\resumeSubSubheading}[2]{
    \item
    \begin{tabular*}{0.97\textwidth}{l@{\extracolsep{\fill}}r}
      \textit{\small#1} & \textit{\small #2} \\
    \end{tabular*}\vspace{-7pt}
}

\newcommand{\resumeProjectHeading}[2]{
    \item
    \begin{tabular*}{0.97\textwidth}{l@{\extracolsep{\fill}}r}
      \small#1 & #2 \\
    \end{tabular*}\vspace{-7pt}
}

\newcommand{\resumeSubItem}[1]{\resumeItem{#1}\vspace{-4pt}}

\renewcommand\labelitemii{$\vcenter{\hbox{\tiny$\bullet$}}$}

\newcommand{\resumeSubHeadingListStart}{\begin{itemize}[leftmargin=0.15in, label={}]}
\newcommand{\resumeSubHeadingListEnd}{\end{itemize}}
\newcommand{\resumeItemListStart}{\begin{itemize}}
\newcommand{\resumeItemListEnd}{\end{itemize}\vspace{-5pt}}

% Role context under a \resumeSubheading. The heading table is top-aligned and
% pulls the next line up, so the blurb must start a new paragraph below it.
\newcommand{\resumeIntro}[1]{\par\vspace{3pt}{\small #1\par}}
.
% Author : Jake Gutierrez
% Based off of: https://github.com/sb2nov/resume
% License : MIT
% This file is not a document. Compile marti_ferran_gonzales_cv.tex or marti_ferran_gonzales_cover_letter.tex.

\ifdefined\cvsetupfont
  \usepackage{fontspec}
  \cvsetupfont
  \renewcommand{\familydefault}{\sfdefault}
\else
  \usepackage[T1]{fontenc}
  \usepackage[utf8]{inputenc}
\fi
\usepackage{latexsym}
\usepackage[empty]{fullpage}
\usepackage{titlesec}
\usepackage{marvosym}
\usepackage[usenames,dvipsnames]{color}
\usepackage{verbatim}
\usepackage{enumitem}
\usepackage[hidelinks]{hyperref}
\usepackage{fancyhdr}
\usepackage[english]{babel}
\usepackage{tabularx}
\ifdefined\cvsetupfont
\else
  \input{glyphtounicode}
\fi

\pagestyle{fancy}
\fancyhf{} % clear all header and footer fields
\fancyfoot{}
\renewcommand{\headrulewidth}{0pt}
\renewcommand{\footrulewidth}{0pt}

% Adjust margins
\addtolength{\oddsidemargin}{-0.5in}
\addtolength{\evensidemargin}{-0.5in}
\addtolength{\textwidth}{1in}
\addtolength{\topmargin}{-.5in}
\addtolength{\textheight}{1.0in}

\urlstyle{same}

\raggedbottom
\raggedright
\setlength{\tabcolsep}{0in}

% Sans faces here have no small caps, so the section title stays in the same font.
\ifdefined\cvsetupfont
  \titleformat{\section}{
    \vspace{-4pt}\raggedright\large
  }{}{0em}{\MakeUppercase}[\color{black}\titlerule \vspace{-5pt}]
\else
  \titleformat{\section}{
    \vspace{-4pt}\scshape\raggedright\large
  }{}{0em}{}[\color{black}\titlerule \vspace{-5pt}]
  % Ensure that the generated pdf is machine readable/ATS parsable
  \pdfgentounicode=1
\fi

%-------------------------
% Custom commands
\newcommand{\resumeItem}[1]{
  \item\small{
    {#1 \vspace{-2pt}}
  }
}

\newcommand{\resumeSubheading}[4]{
  \vspace{-2pt}\item
    \begin{tabular*}{0.97\textwidth}[t]{l@{\extracolsep{\fill}}r}
      \textbf{#1} & #2 \\
      \textit{\small#3} & \textit{\small #4} \\
    \end{tabular*}\vspace{-7pt}
}

\newcommand{\resumeSubSubheading}[2]{
    \item
    \begin{tabular*}{0.97\textwidth}{l@{\extracolsep{\fill}}r}
      \textit{\small#1} & \textit{\small #2} \\
    \end{tabular*}\vspace{-7pt}
}

\newcommand{\resumeProjectHeading}[2]{
    \item
    \begin{tabular*}{0.97\textwidth}{l@{\extracolsep{\fill}}r}
      \small#1 & #2 \\
    \end{tabular*}\vspace{-7pt}
}

\newcommand{\resumeSubItem}[1]{\resumeItem{#1}\vspace{-4pt}}

\renewcommand\labelitemii{$\vcenter{\hbox{\tiny$\bullet$}}$}

\newcommand{\resumeSubHeadingListStart}{\begin{itemize}[leftmargin=0.15in, label={}]}
\newcommand{\resumeSubHeadingListEnd}{\end{itemize}}
\newcommand{\resumeItemListStart}{\begin{itemize}}
\newcommand{\resumeItemListEnd}{\end{itemize}\vspace{-5pt}}

% Role context under a \resumeSubheading. The heading table is top-aligned and
% pulls the next line up, so the blurb must start a new paragraph below it.
\newcommand{\resumeIntro}[1]{\par\vspace{3pt}{\small #1\par}}
